\documentclass[11pt]{article}

\usepackage[margin=1.05in]{geometry}
\usepackage{amsmath}
\usepackage{booktabs}
\usepackage{microtype}
\usepackage[hidelinks]{hyperref}

\title{Discovering a Derivative-Free Optimizer by Evolution\\
Against a Memorisation-Proof Real-World Benchmark}
\author{Peter Cotton}
\date{\today}

\begin{document}
\maketitle

\begin{abstract}
We \emph{discover} rather than hand-design a derivative-free optimizer by evolving a
parameterised DE/ES-hybrid template against a memorisation-proof suite of real-world
objectives. A genetic algorithm searching $14$ behavioural knobs reaches a
normalised regret of $0.110$ versus a Nelder--Mead/Differential-Evolution/CMA-ES
panel (mean rank $1.47$ of $4$), against $0.191$ for the original $12$-knob template.
The gain comes from one added mechanism --- a cheap, R\textsuperscript{2}-gated
separable-quadratic trust-region jump --- which we validate three ways: an A/B
re-evolution, a single-gene toggle ablation (disabling it more than doubles regret),
and independent re-discovery by the production search, which drives the jump's firing
probability to $\approx\!1$. The discovered optimizer is a hedged hybrid, found by
grounding fitness in \emph{real} (disguised) evaluation rather than a synthetic proxy.
\end{abstract}

\section{The question}

Can a competitive derivative-free optimizer be \emph{discovered} by search, rather
than hand-designed --- and discovered against a benchmark it cannot game? We
parameterise a single unified DE/ES-hybrid template and let a genetic algorithm tune
it, scoring every candidate by normalised regret against a fixed baseline panel
(Nelder--Mead, Differential Evolution, CMA-ES) on \emph{disguised} real-world
objectives.

\section{The space of optimizers (the genome)}

A genome is a point in $[0,1]^{14}$ decoding to a concrete optimizer that blends, in
tunable proportion: DE moves (\texttt{rand/1} and current-to-best, mixed),
SHADE-style self-adaptation of $F$/$CR$, local moves (Gaussian or coordinate pattern
search with a $1/5$-success step), simulated-annealing acceptance, restart-on-
stagnation, and --- the two genes added here --- a \textbf{separable-quadratic
trust-region jump}: with probability \texttt{p\_surrogate} a generation fits a
diagonal quadratic $f(x)\approx c+\sum_i b_i x_i+\sum_i a_i x_i^2$ to the nearest
$\sim\!(4n{+}2)$ evaluated points and jumps to its per-coordinate vertex
$x_i^\star=-b_i/(2a_i)$ (only where $a_i>0$), within an adaptive trust radius, gated
by a minimum model $R^2$ (\texttt{r2\_min}). The fit is $O(n)$ --- a poor-man's
NEWUOA step the pure DE/ES moves lack.

\section{The discovery loop}

A $(\mu{+}\lambda)$ genetic algorithm with gene-wise crossover, Gaussian mutation,
elitist survival, and a warm-started, diverse initial population searches the genome
space. Fitness is \emph{continuous} normalised regret (not integer rank) so the
search gets a gradient even when the candidate is merely close behind the panel. The
benchmark is the \texttt{example\_applications} suite, each objective relocated by a
seeded cube$\to$cube diffeomorphism so the optimum cannot be memorised.

\section{Result: the surrogate earns its place}

\begin{table}[h]
\centering
\begin{tabular}{lcc}
\toprule
optimizer & normalised regret & mean rank (of 4) \\
\midrule
base template GA ($12$ knobs) & $0.191$ & $1.73$ \\
surrogate prototype (A/B re-evolution) & $0.112$ & --- \\
\textbf{canonical $14$-knob GA} & $\mathbf{0.110}$ & $\mathbf{1.47}$ \\
\bottomrule
\end{tabular}
\caption{Same scale throughout ($15$ disguised demos, $2$ seeds, $100$ trials).
Lower regret is better; mean rank is among \{candidate, NM, DE, CMA\}.}
\label{tab:results}
\end{table}

Adding the surrogate lowers regret from $0.191$ to $0.110$ (Table~\ref{tab:results}).
Three independent checks attribute the gain to the mechanism itself, not to extra
search:
\begin{enumerate}
\item \textbf{Single-gene ablation.} On the evolved genome, toggling only
\texttt{p\_surrogate}: at $0.625$ (evolved) regret is $0.112$; at $0.0$ (jump off) it
jumps to $0.238$; at $1.0$ (always) it is $0.129$. Disabling the jump \emph{more than
doubles} regret --- and leaves the genome \emph{worse than the base template}
($0.238$ vs $0.191$), showing the base knobs \emph{co-adapted} to rely on the jump
rather than treating it as a bolt-on.
\item \textbf{Independent re-discovery.} The production search, run fresh on the
$14$-knob genome, reaches $0.110$ and drives \texttt{p\_surrogate} to $\approx\!0.998$
--- given free rein, evolution makes the jump fire almost every generation. (It found
a different co-adapted basin than the prototype's $0.625$; both score $\approx\!0.11$,
so the ``optimal'' firing rate is not a constant --- it co-adapts with the base knobs.)
\item \textbf{Behaviour preserved.} The mechanism was folded into the shipped
template additively; the original API and its tests are unchanged.
\end{enumerate}

\section{Caveats}

\begin{itemize}
\item \textbf{In-sample.} Fitness is measured on the same $15$ demos and $2$ seeds the
GA trains on, so the winner could overfit those instances. A train/test split (evolve
on demo-set A, score on disjoint set B and unseen seeds) is running; generalisation is
not yet established.
\item \textbf{Regime-specific.} This is low-budget DFO on \emph{real-world} objectives.
The same curvature-exploiting idea is neutral-to-harmful on
smooth \emph{synthetic} functions (where production CMA dominates), so the surrogate's
value is conditional on the messy, structured, low-budget regime, not universal.
\item \textbf{Honest framing.} What evolution yields is a strong tuned configuration of
an existing hybrid template plus one genuinely new mechanism --- a new optimizer
\emph{instance}, not a new algorithm \emph{family}.
\end{itemize}

\end{document}
